\documentclass[12pt]{article}
\usepackage{hyperref,amsmath,amsfonts,fullpage,color}
\newcommand{\be}{\begin{equation}}
\newcommand{\ee}{\end{equation}}
\newcommand{\bh}{\bar{h}}
\usepackage{enumitem}
\newcounter{quesno}
\newcommand{\ques}[1]{\vspace{3em}\stepcounter{quesno}
\noindent{\bf Question \thequesno.} \emph{(#1 points)\\}
}
\begin{document}

\begin{center}

{\bf\Large General Relativity}

\vspace{1em}
{\bf\large Final Examination}
\vspace{1em}
{\bf Time:} 4 hours


\end{center} 

\vspace{1em}
\hrule
\vspace{1em}
\noindent
{\bf Instructions:} 
For full marks, all questions must be solved.

\vspace{1em}
\hrule

\ques{2} 
The components of the Riemann tensor in some coordinates are $R^\mu{}_{\nu\rho\sigma}$.
Write down how they transform under a change of coordinates.
Define the Ricci tensor by $R_{\mu\nu} = R^\rho{}_{\mu\rho\nu}$, and show explicitly that this transforms as a tensor under change of coordinates.

\ques{3}
Consider a manifold on which we have both a metric, $g_{\mu\nu}$, and a totally antisymmetric $(0,3)$ tensor, $H_{\mu\nu\rho} = H_{[\mu\nu\rho]}$.
We can define a connection with components $\hat\Gamma_{\mu\nu}{}^\rho$ by:
\be
\hat \Gamma_{\mu\nu}{}^\rho = \Gamma_{\mu\nu}{}^\rho + \frac{1}{12} g^{\rho \sigma} H_{\mu\nu\sigma}
\label{stringconn}
\ee
where $\Gamma_{\mu\nu}{}^\rho$ denotes the components of the usual Levi-Civita connection.

\begin{enumerate}[label=\alph*)]
\item (\emph{2 points}) Is the new connection $\hat\Gamma_{\mu\nu}{}^\rho$ defined in equation \eqref{stringconn} torsion-free? Is it metric compatible? 
\item (\emph{1 point}) Show that the geodesics of the connection $\hat\Gamma_{\mu\nu}{}^\rho$ are the same as the geodesics of the Levi-Civita connection. 
\end{enumerate}


\ques{4}
The Schwarzschild metric is:
\be
ds^2 = - f(r) dt^2 + f(r)^{-1} dr^2 + r^2 d \theta^2 + r^2 \sin^2 \theta d\phi^2 
\ee
where $f(r) = 1 - \frac{2GM}{r}$. 

\begin{enumerate}[label=\alph*)]
\item  (\emph{2 points}) Suppose you are sufficiently interested in physics that you decide to fall into a Schwarzschild black hole to see what happens. 
Your friends and family think this might be unsafe, as they have noticed there are singularities in the Schwarzschild metric. 
Where do these singularities occur, which should you be genuinely concerned about, and why? 
\item (\emph{2 points}) Determine the Schwarzschild metric in the ingoing Eddington-Finkelstein coordinates $(v,r,\theta,\phi)$, where
\be
v = t + r + 2GM \ln\left| \frac{r}{2GM} - 1 \right|
\ee
and explain how this supports your answer to part a). 
\end{enumerate} 


\ques{5}
The metric of de Sitter is:
\be
ds^2 =- dt^2 +  L^2  \cosh^2 (t/L) \left( d\chi^2 + \sin^2 \chi \left( d\theta^2 + \sin^2\theta d\phi^2 \right)\right)  \,.
\label{dsmetric} 
\ee
where $t \in (-\infty, \infty)$, $\chi \in (0,\pi)$ and $(\theta,\phi)$ are the usual coordinates on a two-sphere (and together $\chi, \theta,\phi$ parametrise a three-sphere).

\begin{enumerate}[label=\alph*)]
\item (\emph{2 points})  Explain why this metric is a candidate cosmological solution that could describe our universe. What is the scale factor?  
\item (\emph{1 point}) Explain why two metrics that are related by $g_{\mu\nu} = \Omega^2 \bar g_{\mu\nu}$, can be said to have the same causal structure. 
\item (\emph{2 points}) For de Sitter, make the coordinate transformation
\be
\cosh \frac{t}{L} = \frac{1}{\cos \eta} \,,
\ee
and draw the Penrose diagram. 

\end{enumerate}

\ques{6}
The linearised Einstein equation is
\be
- \frac{1}{2} \partial^\rho\partial_\rho \bh_{\mu\nu} + \partial^\rho \partial_{(\mu} \bh_{\nu)\rho}
- \frac{1}{2} \eta_{\mu\nu} \partial^\rho \partial^\sigma \bh_{\rho\sigma} = 8 \pi G T_{\mu\nu} \,,
\ee
when written in terms of $\bh_{\mu\nu} = h_{\mu\nu} - \frac{1}{2} h \eta_{\mu\nu}$. 

\begin{enumerate}[label=\alph*)]
\item (\emph{1 point}) By making an appropriate gauge choice, explain why the linearised Einstein equation leads to \emph{gravitational waves}.
\item (\emph{2 point}) Consider an experiment involving two freely falling test particles, arranged so that one is at spatial position $x^i = (0,0,0)$ and the other at spatial position $x^i=(L_x,0,0)$.
Suppose a gravitational wave travelling along the $z$-axis passes through the laboratory, corresponding to the following metric perturbation:
\be
h_{\mu\nu} = \begin{pmatrix}
0 & 0 & 0 & 0 \\
0 & h_{+}(t-z) & 0 & 0 \\
0 & 0 & - h_{+}(t-z) & 0 \\
0 & 0 & 0 & 0
\end{pmatrix} \,.
\ee
Show that the proper distance between the test particles is given by: 
\be
\Delta s \approx L_x \left( 1 + \frac{1}{2} h_{+}(t) \right) \,.
\ee
If a third test particle is introduced at $x^i=(0,L_y,0)$ show that the proper distance between it and the particle at the origin is
\be
\Delta s \approx L_y \left( 1 - \frac{1}{2} h_{+}(t) \right) \,.
\ee
\item (\emph{3 points}) Consider a binary system consisting of two massive bodies in a circular orbit, which are emitting gravitational waves.
Explain \emph{qualitatively} what happens to the system, and sketch the corresponding gravitational wave signal which would be observed by a detector such as LIGO.
You should mention the following terms: \emph{inspiral}, \emph{merger}, \emph{ringdown}, explain how the interferometer detects a wave such as that described in part b), and connect the features of the signal to the physics of the binary pair. 

\end{enumerate}


\end{document}
